% GENERATED FILE. Edit canonical lessons, then run npm run generate:books.
% canonical-source-hash: fnv1a64:96bad7efee535844
% canonical-lessons: UR-C306-dhamaka, UR-C306-bam, UR-C306-gunjna, UR-C306-tharrana, UR-C306-sadma

\chapter{A Blast}
\label{ch:ur-a-blast}
\hlchaptermodality{306}

\begin{chapteropening}
\textbf{By the end of this chapter:} \emph{I can describe a sudden loud event.}

Complete the last lesson of chapter 306: I can describe a sudden loud event.
\end{chapteropening}

\section[\texorpdfstring{\ur{دھماکا} (dhamākā)}{دھماکا (dhamākā)}]{\ur{دھماکا} (dhamākā) — a bang; a blast}

\label{lesson:UR-C306-dhamaka}



\begin{quote}
Before the new one: say the Urdu for to be startled, then the Urdu for a revelation.
\end{quote}

\subsection*{You'll want to know: \ur{دھماکا}}
\textbf{\ur{دھماکا}} — \emph{dhamākā} — \textquotedblleft{}a bang; a blast\textquotedblright{}.

\begin{quote}
\ur{دھماکا} \ur{ذرا} \ur{دور} \ur{ہوا۔}
\end{quote}

\begin{quote}
\emph{dhamākā zarā dūr huā.}
\end{quote}

\begin{quote}
The blast happened a little way off.
\end{quote}

Masculine.

\begin{cousinweb}
\textbf{\ur{دھماکا}} (\emph{dhamākā}) is probably imitative, built on the sound \emph{dham}; CDIAL sets Hindi \emph{dhamādham}, \textquotedblleft{}noisily\textquotedblright{}, under a form it reconstructs (CDIAL 6735).
\end{cousinweb}

\begin{grammarlens}[title={What you've built}]
One more word for telling what happened.
\end{grammarlens}

\subsection*{Your turn}
\begin{itemize}
\raggedright
  \item \emph{Say it:} \emph{dhamākā}
  \item \emph{Say it:} \emph{dhamākā}, once more
  \item \emph{Say it:} \emph{dhamākā zarā dūr huā.} — The blast happened a little way off.
  \item \emph{Recall:} say the Urdu for to be startled, then the Urdu for a revelation, then say \emph{dhamākā} again
\end{itemize}

\begin{tcolorbox}[breakable,colback=teal!4,colframe=teal!35!black,title={Before you move on}]
What does \ur{دھماکا} mean? (\textquotedblleft{}A bang; a blast\textquotedblright{}.) Say it once more.
\end{tcolorbox}
\section[\texorpdfstring{\ur{بم} (bam)}{بم (bam)}]{\ur{بم} (bam) — a bomb}

\label{lesson:UR-C306-bam}



\begin{quote}
Before the new one: say the Urdu for a revelation, then the Urdu for a blast.
\end{quote}

\subsection*{You'll want to know: \ur{بم}}
\textbf{\ur{بم}} — \emph{bam} — \textquotedblleft{}a bomb\textquotedblright{}.

\begin{quote}
\ur{پولیس} \ur{کو} \ur{ایک} \ur{بم} \ur{ملا۔}
\end{quote}

\begin{quote}
\emph{polīs ko ek bam milā.}
\end{quote}

\begin{quote}
The police found a bomb.
\end{quote}

Masculine.

\begin{cousinweb}
\textbf{\ur{بم}} (\emph{bam}) is borrowed from English \emph{bomb}. Urdu says it with a short \emph{a}.
\end{cousinweb}

\begin{grammarlens}[title={What you've built}]
One more word for telling what happened.
\end{grammarlens}

\subsection*{Your turn}
\begin{itemize}
\raggedright
  \item \emph{Say it:} \emph{bam}
  \item \emph{Say it:} \emph{bam}, once more
  \item \emph{Say it:} \emph{polīs ko ek bam milā.} — The police found a bomb.
  \item \emph{Recall:} say the Urdu for a revelation, then the Urdu for a blast, then say \emph{bam} again
\end{itemize}

\begin{tcolorbox}[breakable,colback=teal!4,colframe=teal!35!black,title={Before you move on}]
What does \ur{بم} mean? (\textquotedblleft{}A bomb\textquotedblright{}.) Say it once more.
\end{tcolorbox}
\section[\texorpdfstring{\ur{گونجنا} (gūnjnā)}{گونجنا (gūnjnā)}]{\ur{گونجنا} (gūnjnā) — to echo; to resound}

\label{lesson:UR-C306-gunjna}



\begin{quote}
Before the new one: say the Urdu for a blast, then the Urdu for a bomb.
\end{quote}

\subsection*{You'll want to know: \ur{گونجنا}}
\textbf{\ur{گونجنا}} — \emph{gūnjnā} — \textquotedblleft{}to echo; to resound\textquotedblright{}.

\begin{quote}
\ur{دھماکے} \ur{سے} \ur{گلی} \ur{گونج} \ur{اٹھی۔}
\end{quote}

\begin{quote}
\emph{dhamāke se galī gūnj uṭhī.}
\end{quote}

\begin{quote}
The lane rang with the blast.
\end{quote}

Past (it …): \ur{وہ} \ur{گلی} \ur{گونجی}, \emph{voh galī gūnjī}.

\begin{cousinweb}
\textbf{\ur{گونجنا}} (\emph{gūnjnā}) is inherited: CDIAL lists the Hindi cousin \emph{gū̃jnā}, \textquotedblleft{}to hum\textquotedblright{}, under Sanskrit \emph{gúñjati} (CDIAL 4175); in Urdu a sound that fills a place \textbf{\ur{گونجتی} \ur{ہے}}.
\end{cousinweb}

\begin{grammarlens}[title={What you've built}]
One more word for telling what happened.
\end{grammarlens}

\subsection*{Your turn}
\begin{itemize}
\raggedright
  \item \emph{Say it:} \emph{gūnjnā}
  \item \emph{Say it:} \emph{gūnjnā}, once more
  \item \emph{Say it:} \emph{dhamāke se galī gūnj uṭhī.} — The lane rang with the blast.
  \item \emph{Recall:} say the Urdu for a blast, then the Urdu for a bomb, then say \emph{gūnjnā} again
\end{itemize}

\begin{tcolorbox}[breakable,colback=teal!4,colframe=teal!35!black,title={Before you move on}]
What does \ur{گونجنا} mean? (\textquotedblleft{}To echo; to resound\textquotedblright{}.) Say it once more.
\end{tcolorbox}
\section[\texorpdfstring{\ur{تھرانا} (tharrānā)}{تھرانا (tharrānā)}]{\ur{تھرانا} (tharrānā) — to quake; to shudder}

\label{lesson:UR-C306-tharrana}



\begin{quote}
Before the new one: say the Urdu for a bomb, then the Urdu for to echo.
\end{quote}

\subsection*{You'll want to know: \ur{تھرانا}}
\textbf{\ur{تھرانا}} — \emph{tharrānā} — \textquotedblleft{}to quake; to shudder\textquotedblright{}.

\begin{quote}
\ur{دھماکے} \ur{سے} \ur{کھڑکیاں} \ur{تھرا} \ur{گئیں۔}
\end{quote}

\begin{quote}
\emph{dhamāke se khiṛkiyāṅ tharrā gaīṅ.}
\end{quote}

\begin{quote}
The windows shook with the blast.
\end{quote}

Past (he …): \ur{وہ} \ur{تھرایا}, \emph{voh tharrāyā}.

\begin{cousinweb}
\textbf{\ur{تھرانا}} (\emph{tharrānā}) is inherited: CDIAL lists the Hindi cousin \emph{tharrānā}, \textquotedblleft{}to tremble\textquotedblright{}, under a form it reconstructs (CDIAL 6092).
\end{cousinweb}

\begin{grammarlens}[title={What you've built}]
One more word for telling what happened.
\end{grammarlens}

\subsection*{Your turn}
\begin{itemize}
\raggedright
  \item \emph{Say it:} \emph{tharrānā}
  \item \emph{Say it:} \emph{tharrānā}, once more
  \item \emph{Say it:} \emph{dhamāke se khiṛkiyāṅ tharrā gaīṅ.} — The windows shook with the blast.
  \item \emph{Recall:} say the Urdu for a bomb, then the Urdu for to echo, then say \emph{tharrānā} again
\end{itemize}

\begin{tcolorbox}[breakable,colback=teal!4,colframe=teal!35!black,title={Before you move on}]
What does \ur{تھرانا} mean? (\textquotedblleft{}To quake; to shudder\textquotedblright{}.) Say it once more.
\end{tcolorbox}
\section[\texorpdfstring{\ur{صدمہ} (sadmā)}{صدمہ (sadmā)}]{\ur{صدمہ} (sadmā) — a shock (of grief or fright)}

\label{lesson:UR-C306-sadma}



\begin{quote}
Before the new one: say the Urdu for to echo, then the Urdu for to shudder.
\end{quote}

\subsection*{You'll want to know: \ur{صدمہ}}
\textbf{\ur{صدمہ}} — \emph{sadmā} — \textquotedblleft{}a shock (of grief or fright)\textquotedblright{}.

\begin{quote}
\ur{ماں} \ur{کو} \ur{خبر} \ur{سے} \ur{صدمہ} \ur{ہوا۔}
\end{quote}

\begin{quote}
\emph{māṅ ko khabar se sadmā huā.}
\end{quote}

\begin{quote}
The news came as a shock to the mother.
\end{quote}

Masculine.

\begin{cousinweb}
\textbf{\ur{صدمہ}} (\emph{sadmā}) is usually traced, through Persian, to Arabic \emph{ṣadma}, \textquotedblleft{}a blow, a collision\textquotedblright{}; in Urdu it is the blow of bad news. The one shocked takes \textbf{\ur{کو}}.
\end{cousinweb}

\begin{grammarlens}[title={What you've built}]
One more word for telling what happened.
\end{grammarlens}

\subsection*{Your turn}
\begin{itemize}
\raggedright
  \item \emph{Say it:} \emph{sadmā}
  \item \emph{Say it:} \emph{sadmā}, once more
  \item \emph{Say it:} \emph{māṅ ko khabar se sadmā huā.} — The news came as a shock to the mother.
  \item \emph{Recall:} say the Urdu for to echo, then the Urdu for to shudder, then say \emph{sadmā} again
\end{itemize}

\begin{tcolorbox}[breakable,colback=teal!4,colframe=teal!35!black,title={Before you move on}]
What does \ur{صدمہ} mean? (\textquotedblleft{}A shock (of grief or fright)\textquotedblright{}.) Say it once more.
\end{tcolorbox}
